\subsection{预训练机制：预测下一个 token 为何能够泛化}

从 slide 7 开始，课程把“预测下一个词”拆成可执行过程。自回归语言模型先把文本转换成 token，做一次前向计算得到下一个 token 的条件分布，再采样并解码。n-gram 是最直接的统计基线：统计上下文后各种续词的频率。但自然语言组合近乎无限，绝大多数长上下文只出现一次，显式计数既占空间又无法迁移到未见句子。神经语言模型的关键贡献，是把离散上下文压缩到连续表示中，让相似语义共享统计强度。

\lecturefigure{slides-images/slide-007.png}{官方 slide 7：预训练目标、数据规模与历史位置}
\lecturefigure{slides-images/slide-008.png}{官方 slide 8：自回归语言模型的 tokenization、forward、sampling 与 detokenization}
\lecturefigure{slides-images/slide-009.png}{官方 slide 9：n-gram 条件频率估计及其泛化失败}
\lecturefigure{slides-images/slide-010.png}{官方 slide 10：神经语言模型作为 n-gram 稀疏性的解决方向}

\begin{importantbox}{从计数到参数共享}
对 n-gram，条件概率可写为
\[
P(x\mid c)=\frac{\operatorname{Count}(c,x)}{\operatorname{Count}(c)},
\]
其中 $c$ 表示前文上下文，$x$ 表示候选下一个 token。问题不在公式错误，而在每个 $c$ 都需要独立计数。神经模型学习参数化函数 $P_\theta(x\mid c)$，使不同上下文通过 embedding 和网络层共享参数；它不是记住所有句子，而是学习哪些上下文在预测上应被视为相似。
\end{importantbox}

\begin{knowledgebox}{背景概念：token、word 与 character}
token 是 tokenizer 输出的离散符号，不等于自然语言中的“词”。英文常见 token 约为若干字符，中文可能是单字、词片段或字节组合。训练规模用 token 计，是因为模型每一步预测的是词表中的 token ID；同一段文字在不同 tokenizer 下会产生不同长度，所以“10T tokens”不能无条件换算成固定字数。slide 8 的五步流程也说明，tokenization 属于模型输入协议的一部分，不是纯粹的预处理细节。
\end{knowledgebox}

\textbf{老师强调：}讲者把预训练描述为“尽可能学习互联网和世界”，这是训练直觉，不是知识完备性保证。next-token objective 会同时学习事实、语言结构、偏见和数据中的错误；模型能流畅续写，也不意味着它已经具备可靠的用户意图理解或行动约束。这正是后训练存在的原因，下一组 slides 会先说明预训练语料如何被构造，再讨论目标如何改变。

